\section{表示}

\subsection{表示}

\begin{definition}[表示]\label{def:representation}
设 $A$ 为 $\mathbb{F}$-代数. $A$ 的一个表示是一个二元组 $(V,\rho)$, 其中
\begin{enumerate}
  \item $V$ 是 $\mathbb{F}$-线性空间;
  \item $\rho\colon A\to\operatorname{End}(V)$, $a\mapsto\rho(a)$ 是代数同态.
\end{enumerate}
有时简称 $V$ 为 $A$ 的一个表示, 或称 $\rho$ 为 $A$ 的一个表示.
\end{definition}

\begin{example}\label{ex:representation-FZ2}
\mbox{}\par\vspace{-0.2em}
\begin{enumerate}
  \item 由 $\mathbb{F}\mathbb{Z}_2\to M_2(\mathbb{F})\cong\operatorname{End}(\mathbb{F}^2)$,
  \[
    \rho\colon g\longmapsto\begin{pmatrix}0&1\\1&0\end{pmatrix},
  \]
  得到 $\mathbb{F}\mathbb{Z}_2$ 的一个表示 $(\mathbb{F}^2,\rho)$.
  \item 由 $\mathbb{C}\mathbb{Z}_n\to\operatorname{End}(\mathbb{C})=\mathbb{C}$,
  \[
    \rho\colon g^j\longmapsto\xi^j,\qquad \xi=\exp\left(\frac{2\pi i}{n}\right),
  \]
  得到 $\mathbb{C}\mathbb{Z}_n$ 的一个表示 $(\mathbb{C},\rho)$.
\end{enumerate}
\end{example}

\begin{theorem}\label{thm:rep-is-module}
$V$ 是 $A$ 的一个表示当且仅当 $V$ 是一个 $A$-模.
\end{theorem}

\begin{proof}
\mbox{}\par\vspace{-0.2em}
\begin{enumerate}
\item 设 $V$ 为 $A$ 的一个表示, 有同态 $\rho\colon A\to\operatorname{End}(V)$, $a\mapsto\rho(a)$.
定义运算
\[
  A\times V\longrightarrow V,\qquad (a,v)\longmapsto a\cdot v:=\rho(a)(v).
\]
于是
\begin{align*}
  (a_1+a_2)\cdot v&=\rho(a_1+a_2)(v)=(\rho(a_1)+\rho(a_2))(v)\\
  &=\rho(a_1)(v)+\rho(a_2)(v)=a_1\cdot v+a_2\cdot v,\\
  (a_1a_2)\cdot v&=\rho(a_1a_2)(v)=(\rho(a_1)\rho(a_2))(v)\\
  &=\rho(a_1)\bigl(\rho(a_2)(v)\bigr)=a_1\cdot(a_2\cdot v),\\
  (ka)\cdot v&=\rho(ka)(v)=k\rho(a)(v)=k(a\cdot v),\\
  1_A\cdot v&=\rho(1_A)(v)=\operatorname{Id}(v)=v,\\
  a\cdot(v_1+v_2)&=\rho(a)(v_1+v_2)=\rho(a)(v_1)+\rho(a)(v_2)\\
  &=a\cdot v_1+a\cdot v_2,\\
  a\cdot(kv)&=\rho(a)(kv)=k\rho(a)(v)=k(a\cdot v).
\end{align*}
故 $V$ 是一个 $A$-模.

\item 设 $V$ 为 $A$-模, 运算为 $A\times V\to V$, $(a,v)\mapsto a\cdot v$.
定义映射 $\rho\colon A\to\operatorname{End}(V)$, $a\mapsto\rho(a)$, 其中 $\rho(a)$ 由 $\rho(a)(v)=a\cdot v$ 规定.
只需证明 $\rho$ 是代数同态.
\begin{enumerate}
  \item $\rho$ 是线性映射:
  $\rho(a_1+a_2)(v)=(a_1+a_2)\cdot v=a_1\cdot v+a_2\cdot v=\rho(a_1)(v)+\rho(a_2)(v)$,
  即 $\rho(a_1+a_2)=\rho(a_1)+\rho(a_2)$;
  $\rho(ka)(v)=(ka)\cdot v=k(a\cdot v)=k\rho(a)(v)$, 即 $\rho(ka)=k\rho(a)$.
  \item $\rho$ 保持乘法:
  $\rho(ab)(v)=(ab)\cdot v=a\cdot(b\cdot v)=\rho(a)\bigl(\rho(b)(v)\bigr)=(\rho(a)\rho(b))(v)$,
  即 $\rho(ab)=\rho(a)\rho(b)$.
\end{enumerate}
故 $V$ 是 $A$ 的表示.
\end{enumerate}
\end{proof}

\begin{definition}[群的表示]\label{def:group-representation}
设 $G$ 为群, $V$ 为 $\mathbb{F}$-线性空间. 称 $V$ 为 $G$ 的一个 (线性) 表示, 若存在代数同态
\[
  \rho\colon KG\longrightarrow\operatorname{End}(V).
\]
\end{definition}

\begin{example}\label{ex:FZ2-representation-model}
\mbox{}\par\vspace{-0.2em}
\begin{enumerate}
  \item 设 $\rho\colon \mathbb{F}\mathbb{Z}_2\to\operatorname{End}(\mathbb{F}^2)\cong M_2(\mathbb{F})$, 其中
  $\mathbb{Z}_2=\langle g\mid g^2=1\rangle$, $\mathbb{F}^2=\operatorname{span}_{\mathbb{F}}\{\alpha_1,\alpha_2\}$,
  $\alpha_1=(1,0)^{\mathrm{T}}$, $\alpha_2=(0,1)^{\mathrm{T}}$, 且
  \[
    \rho(g)=\begin{pmatrix}0&1\\1&0\end{pmatrix}.
  \]
  规定 $g\cdot v=\rho(g)(v)$, 则
  \begin{align*}
    g\cdot\alpha_1&=\rho(g)(\alpha_1)=\begin{pmatrix}0&1\\1&0\end{pmatrix}\binom10=\binom01=\alpha_2,\\
    g\cdot\alpha_2&=\rho(g)(\alpha_2)=\begin{pmatrix}0&1\\1&0\end{pmatrix}\binom01=\binom10=\alpha_1 .
  \end{align*}
  \item 正则模: $\mathbb{F}\mathbb{Z}_2\curvearrowright \mathbb{F}\mathbb{Z}_2=\operatorname{span}_{\mathbb{F}}\{1,g\}$, 记 $\mathbb{F}\mathbb{Z}_2=V$,
  \[
    \rho\colon \mathbb{F}\mathbb{Z}_2\longrightarrow\operatorname{End}(V),\qquad g\longmapsto\rho(g).
  \]
  考察 $\rho(g)$ 对基 $\{1,g\}$ 的作用: $1\mapsto g$, $g\mapsto 1$, 对应的矩阵为
  $\begin{pmatrix}0&1\\1&0\end{pmatrix}$. 可见把 $g$ 映为该矩阵是合理的.
\end{enumerate}
\end{example}

\begin{proposition}\label{prop:regular-module-exists}
任何群代数都有正则模 (都对应唯一同态).
\end{proposition}

\subsection{表示论的直观解释}

以矩阵来说明. 设 $\rho\colon A\to\operatorname{End}(V)$ 为表示, 取子空间 $V_1\leq V$.
设 $V_1$ 的一组基为 $\alpha_1,\alpha_2,\dots,\alpha_m$, 那么可以扩充为 $V$ 的一组基
$\alpha_1,\dots,\alpha_m,\alpha_{m+1},\dots,\alpha_n$.

\begin{definition}[子表示]\label{def:subrepresentation}
称 $\rho|_{V_1}$ 为 $\rho$ 的子表示, 若对任意 $a\in A$, $v_1\in V_1$, 都有 $\rho(a)(v_1)\in V_1$.
\end{definition}

此时 $V_1$ 是 $\rho$-不变子空间, 且
\[
  \rho(a)(\alpha_1,\dots,\alpha_m,\alpha_{m+1},\dots,\alpha_n)
  =(\alpha_1,\dots,\alpha_m,\alpha_{m+1},\dots,\alpha_n)
  \begin{pmatrix}
    M_1 & C\\
    0 & M_2
  \end{pmatrix},
\]
其中 $M_1\in M_m(\mathbb{F})$, $M_2\in M_{n-m}(\mathbb{F})$. 由于 $\rho|_{V_1}$ 只对 $\alpha_1,\dots,\alpha_m$ 有作用,
故 $M_1$ 描述了子表示 (在嵌入意义下); 而在商空间 $V/V_1$ 中, $V_1$ 的元素 (由基
$\alpha_1,\dots,\alpha_m$ 给出) 被映为 $0$, 故 $V/V_1$ 可由 $M_2$ 描述 (在嵌入意义下).

若出现零块 $0$, 则必出现子表示 $M_1$. 因此若没有零块 $0$, 就找不到非平凡的子表示,
即该表示不可约.

当左上角 $C=0$ 时, 矩阵变为准对角矩阵
\[
  \begin{pmatrix}
    * & 0\\
    0 & *
  \end{pmatrix},
\]
记 $M_1$ 对应子空间 $V_1$, $M_2$ 对应子空间 $V_2$, 则可写 $V=V_1\oplus V_2$,
且 $V/V_1\cong V_2$, $V/V_2\cong V_1$. 此时称表示可分解; 否则称其不可分解,
此时不会出现准对角的情形.

\begin{example}\label{ex:FZ2-semisimple-decomposition}
设 $\mathbb{F}\mathbb{Z}_2\curvearrowright \mathbb{F}\mathbb{Z}_2=\operatorname{span}_{\mathbb{F}}\{1,g\}$.
先设 $\operatorname{char}\mathbb{F}\neq2$, 记
\[
  e_1=\frac12(1+g),\qquad e_2=\frac12(1-g).
\]
则 $V_1=\operatorname{span}_{\mathbb{F}}\{e_1\}$, $V_2=\operatorname{span}_{\mathbb{F}}\{e_2\}$ 均为子表示, 并且
\begin{align*}
  g\cdot e_1&=\frac12 g\cdot(1+g)=\frac12(g+g^2)=\frac12(1+g)=e_1,\\
  g\cdot e_2&=\frac12 g\cdot(1-g)=\frac12(g-1)=-e_2 .
\end{align*}
于是 $\mathbb{F}\mathbb{Z}_2$ 可以分解为 $\mathbb{F}\mathbb{Z}_2=V_1\oplus V_2$. 既然 $\mathbb{F}\mathbb{Z}_2$ 有真子模,
故 $\mathbb{F}\mathbb{Z}_2$ 可约.

再考虑 $\operatorname{char}\mathbb{F}=2$. 此时 $\mathbb{F}$ 中 $1=-1$, 取 $e_1=1+g$, 有 $g\cdot e_1=e_1$,
记 $V_1=\operatorname{span}_{\mathbb{F}}\{e_1\}$, 这是 $\mathbb{F}\mathbb{Z}_2$ 的子模, 因此 $\mathbb{F}\mathbb{Z}_2$ 同样可约.
但此时 $\mathbb{F}\mathbb{Z}_2$ 不可分解: 否则设 $\mathbb{F}\mathbb{Z}_2=M_1\oplus M_2$, 其中
$\dim M_1=\dim M_2=1$. 设 $M_1=\operatorname{span}_{\mathbb{F}}\{\alpha_1\}$, $M_2=\operatorname{span}_{\mathbb{F}}\{\alpha_2\}$,
并设 $g\cdot\alpha_1=\lambda\alpha_1$, 则 $g^2\cdot\alpha_1=\alpha_1=\lambda^2\alpha_1$, 故 $\lambda^2=1$,
即 $\lambda=1$ (因为 $-1=1$). 同理 $g\cdot\alpha_2=\alpha_2$. 于是 $g$ 在 $\mathbb{F}\mathbb{Z}_2$ 上为恒同映射,
特别地 $g\cdot1=1$, 但 $g\cdot1=g$, 矛盾. 故 $\operatorname{char}\mathbb{F}=2$ 时 $\mathbb{F}\mathbb{Z}_2$ 不可分解.
\end{example}

\subsection{表示同态}

模与表示本质上是一样的, 这里仅给出描述同态的不同语言.
设表示 $\rho_i\colon A\to\operatorname{End}(V_i)$, 对应运算
\[
  A\times V_i\longrightarrow V_i,\qquad (a,v_i)\longmapsto a\cdot v_i:=\rho_i(a)(v_i),
\]
其中 $a\in A$, $v_i\in V_i$. 回忆模同态 $f\colon V_1\to V_2$ 的要求: 涉及模运算的第 2 条, 即
\[
  f(a\cdot m)=a\cdot f(m),\qquad a\in A,\ m\in V_1 .
\]
在表示的语言下可写为
\[
  f\bigl(\rho_1(a)(m)\bigr)=\rho_2(a)\bigl(f(m)\bigr),\qquad a\in A,\ m\in V_1 .
\]

\subsection{忠实表示}

\begin{definition}[忠实表示]\label{def:faithful-rep}
称表示 $\rho\colon A\to\operatorname{End}(V)$ 是忠实的, 若 $\rho$ 是单射.
\end{definition}

\begin{proposition}\label{prop:faithful-equiv}
$V$ 是忠实 $A$-模当且仅当 $V$ 是忠实 $A$-表示.
\end{proposition}

\begin{proof}
\mbox{}\par\vspace{-0.2em}
\begin{enumerate}
\item 设 $V$ 为忠实模, 只要证明 $\rho$ 是单射, 即 $\operatorname{Ker}(\rho)=\{0\}$.
任取 $a\in\operatorname{Ker}(\rho)$, 则 $\rho(a)=0$, 于是 $\rho(a)(v)=0$, $\forall v\in V$,
即 $a\cdot v=0$. 由忠实模的定义知 $a=0$. 故 $\operatorname{Ker}(\rho)=\{0\}$, $\rho$ 是单射,
$V$ 是忠实表示.

\item 设 $V$ 为忠实表示, 且 $a\cdot v=0$, $\forall v\in V$. 由表示的定义, 这意味着
$\rho(a)(v)=0$, $\forall v\in V$, 即 $\rho(a)=0$. 而 $\rho$ 是单射, 故 $a=0$. 因此 $V$ 是忠实模.
\end{enumerate}
\end{proof}

\begin{remark}\label{rem:faithful-algebra-only}
上述命题中特意强调了 "代数", 这是因为在其他结构中结论不一定成立.
例如在群表示中表示是 $\rho\colon G\to\operatorname{End}(V)$, 而对模而言仍定义在群代数 $\mathbb{F}G$ 上;
扩大 $\rho$ 的作用范围, 可能不会保证 $\rho$ 仍是单射.
\end{remark}

\begin{example}\label{ex:regular-module-faithful}
设 $A$ 为有限维代数, 则总有正则模 ${}_AA$. 若 $a\cdot b=0$ 对一切 $b\in A$ 成立,
取 $b=1$ 即得 $a=0$. 故正则模一定是忠实模.
\end{example}

\begin{theorem}\label{thm:algebra-embeds-in-matrix}
任何有限维代数一定是某个全矩阵代数的子代数.
\end{theorem}

\begin{proof}
设 $A$ 为有限维代数, 则 $V={}_AA$ 一定是忠实模, 从而 $V$ 是忠实表示.
于是存在单同态 $\rho\colon A\to\operatorname{End}(V)\cong M_n(\mathbb{F})$. 视 $\rho$ 为嵌入即得证.
\end{proof}
